\documentclass[12pt]{article}
\usepackage{../../ficha}

\begin{document}

\ficha{Juros Compostos e Função Exponencial}

\instrucoes{Resolve os exercícios. Podes usar calculadora nos problemas; arredonda às centésimas. Em regime de juro composto, $C_n = C_0(1 + t)^n$.}

\secao{Juros Compostos}{Resolve.}
\begin{exercicios}
  \item Um capital de 1\,000 euros é aplicado à taxa anual de 4\,\% em juro composto. Qual é o capital ao fim de 1, 2 e 5 anos?
  \item Compara com o capital acumulado em juro simples, nas mesmas condições, ao fim de 5 anos.
  \item Que capital se deve aplicar hoje, à taxa anual de 3\,\%, para ter 10\,000 euros dentro de 6 anos?
  \item Um capital de 2\,000 euros, à taxa anual de 5\,\%, capitalizada mensalmente ($t = \frac{0{,}05}{12}$ por mês), quanto vale ao fim de 3 anos?
  \item Um carro de 25\,000 euros desvaloriza 15\,\% ao ano. Quanto vale ao fim de 4 anos?
  \item A população de uma vila cresce 2\,\% ao ano. Hoje tem 12\,000 habitantes. Quantos terá dentro de 10 anos?
\end{exercicios}

\secao{Potências e Equações Exponenciais Simples}{Resolve em $\R$.}
\begin{exercicios}
  \item $2^x = 32$
  \item $3^{x+1} = 81$
  \item $5^{2x} = \dfrac{1}{25}$
  \item $4^x = 8$
  \item $9^{x} = 27^{x - 1}$
  \item $2^{x^2} = 16$
  \item $e^{3x - 1} = 1$
  \item $\left(\dfrac{1}{2}\right)^x = 4^{x+3}$
  \item $2^{x+2} + 2^x = 40$
  \item $4^x - 6 \cdot 2^x + 8 = 0$
\end{exercicios}

\secao{Inequações Exponenciais}{Resolve em $\R$.}
\begin{exercicios}
  \item $2^x > 8$
  \item $3^{x - 2} \le 9$
  \item $\left(\dfrac{1}{3}\right)^x > 27$
  \item $e^{2x} < e^{x + 3}$
  \item $0{,}5^{x} \ge 0{,}25$
  \item $4^x - 2^x - 2 > 0$
\end{exercicios}

\secao{A Função Exponencial}{Resolve.}
\begin{exercicios}
  \item Indica o domínio, o contradomínio, a assíntota e a monotonia de $f(x) = 2^x$ e de $g(x) = \left(\dfrac{1}{2}\right)^x$.
  \item Esboça, no mesmo referencial, $y = 2^x$, $y = e^x$ e $y = 3^x$. Em que ponto se intersetam?
  \item Indica o domínio, o contradomínio e a assíntota de $f(x) = 3 - 2^{x+1}$.
  \item Determina o zero de $f(x) = e^x - 1$ e o de $g(x) = 2^{x} - 8$.
  \item Determina $a$ e $k$ de modo que o gráfico de $f(x) = k \cdot a^x$ passe nos pontos $(0, 3)$ e $(2, 12)$.
  \item Mostra que $f(x) = \dfrac{e^x + e^{-x}}{2}$ é uma função par.
\end{exercicios}

\secao{Desafios}{Resolve os problemas.}
\begin{exercicios}
  \item Ao fim de quantos anos (completos) um capital duplica à taxa anual de 6\,\% em juro composto? (Usa a calculadora por tentativas.)
  \item Uma substância radioativa perde metade da massa a cada 8 dias. Começando com 200\,g, escreve a massa $m(t)$ ao fim de $t$ dias e calcula a massa ao fim de 20 dias.
  \item Um banco oferece 4\,\% ao ano com capitalização anual; outro oferece 3,9\,\% ao ano com capitalização mensal. Qual é mais vantajoso?
  \item Resolve: $3^{2x} - 4 \cdot 3^{x} + 3 = 0$.
\end{exercicios}

\end{document}
